% Ipcha Mistabra: Structured Adversarial Verification as a Defense
% Against Sycophancy in Multi-Agent LLM Systems
%
% arXiv preprint — compile with pdflatex
% ---------------------------------------------------------------

\documentclass[11pt,a4paper]{article}

% ---- Encoding & Fonts ----
\usepackage[utf8]{inputenc}
\usepackage[T1]{fontenc}
\usepackage{lmodern}

% ---- Page geometry (arXiv default is fine, but explicit is safer) ----
\usepackage[margin=2.5cm]{geometry}

% ---- Math ----
\usepackage{amsmath,amssymb,amsfonts}

% ---- Tables ----
\usepackage{booktabs}
\usepackage{array}
% \usepackage{multirow}  % not used in this paper

% ---- Graphics ----
\usepackage{graphicx}

% ---- Code listings ----
\usepackage{listings}
\lstset{
  basicstyle=\ttfamily\small,
  keywordstyle=\bfseries,
  frame=single,
  breaklines=true,
  columns=fullflexible,
  language=Python
}

% ---- Hyperlinks (load last-ish) ----
\usepackage[colorlinks=true,linkcolor=blue!60!black,citecolor=blue!60!black,urlcolor=blue!60!black]{hyperref}

% ---- Bibliography ----
\usepackage[numbers,sort&compress]{natbib}

% ---- Misc ----
\usepackage{enumitem}
\usepackage{xcolor}
\usepackage{microtype}

% ---- Title block ----
\title{%
  Ipcha Mistabra: Structured Adversarial Verification\\
  as a Defense Against Sycophancy\\
  in Multi-Agent LLM Systems%
}

\author{%
  Oliver Baer\\
  Independent Researcher, nyxCore Project\\
  \texttt{oliver@nyxcore.cloud}%
}

\date{April 2026 (revised)}

% ================================================================
\begin{document}
\maketitle

% ----------------------------------------------------------------
\begin{abstract}
Large Language Models (LLMs) exhibit sycophancy---the systematic
alignment of outputs with perceived user expectations at the expense
of correctness.  Current mitigations (RLHF, constitutional AI,
prompt-level strategies) remain model-internal properties vulnerable
to distributional shift.  We present the \textbf{Ipcha Mistabra
Protocol (IPCHA)}, a structured adversarial verification framework
that elevates sycophancy defense from a model property to an
\emph{architectural constraint}.  IPCHA implements mandatory
adversarial challenge through a three-agent pipeline (Proponent,
Ipcha Agent, Auditor) with model diversity enforcement and authority
grounding.

We define the \textbf{Integrity Score (IS)}, a family of
domain-specific metrics quantifying claim--evidence alignment
within the IPCHA pipeline (distinct from existing integrity
metrics in other contexts, such as Reasoning Integrity
Scores in clinical AI), and present an empirical
upgrade from TF-IDF cosine similarity to Natural Language Inference
(NLI) via a DeBERTa-v3-based microservice, building on established
NLI-based factual consistency methods (AlignScore~\cite{zha2023alignscore},
SummaC~\cite{laban2022summac}) and integrating them into a
structurally adversarial context.  Evaluation on a controlled
synthetic corpus of 110~cases (55 ACCEPTED, 55 REJECTED) across
8~categories demonstrates that NLI-based scoring produces
significantly wider score separation between legitimate and
adversarial claims compared to TF-IDF ($4.1\times$, Wilcoxon
$p<0.001$, Cohen's $d=-0.43$) and produces significantly different
score distributions from sentence-embedding baselines ($p<0.001$,
Cohen's $d=-0.80$, large effect).  While binary classification accuracy is
high for all methods (TF-IDF: 98.18\%, SBERT: 100\%, NLI: 100\%),
the differences are not statistically significant at this sample size
(McNemar $p>0.05$), and the uniformly high accuracy reflects
synthetic corpus simplicity rather than adversarial robustness; the
primary contribution is the score confidence margin, not the accuracy
delta.  IS interpretation bands are empirically derived via quartile
analysis, replacing previously arbitrary thresholds.

We discuss architectural contributions (fail-closed coherence
validation, sycophancy monitoring, denial-of-wallet defense)
alongside honest limitations including unverified model-family
independence assumptions, the oracle problem, NLI performance
degradation on domain-specific technical claims, statistical
underpowering of the evaluation corpus, circular weight calibration,
audit infrastructure requirements, and the fundamental impossibility
of self-referential verification completeness.
\end{abstract}

\smallskip
\noindent\textbf{Keywords:}
sycophancy mitigation,
adversarial verification,
multi-agent LLM,
natural language inference,
claim verification,
epistemic integrity

% ================================================================
\section{Introduction}\label{sec:intro}

\subsection{The Sycophancy Problem}

The deployment of Large Language Models in decision-critical
workflows---from code review to regulatory compliance
auditing---has exposed a fundamental tension between model
helpfulness and epistemic integrity.  \emph{Sycophancy}, defined as
the tendency of an LLM to produce outputs that align with perceived
user preferences rather than ground truth, represents a systematic
failure mode that compounds over time.

Sharma et al.~\cite{sharma2024sycophancy} characterise sycophancy
through controlled experiments demonstrating that models
systematically shift their responses to match user opinions, even
when initially correct.  Their work identifies distinct behavioral
patterns including feedback-driven response modification and
opinion-matching.  These patterns persist across model scales and
alignment approaches.

The danger in production systems is not that the model produces
incorrect outputs---error is expected and manageable---but that it
produces \emph{confidently incorrect outputs indistinguishable from
correct ones}.  In a code review pipeline, sycophantic agreement with
a flawed architectural proposal can propagate through downstream
steps, encoding the original flaw as an established pattern.

\subsection{Limitations of Current Approaches}

Existing sycophancy mitigations operate at three levels, each with
fundamental limitations:

\textbf{Model-level alignment}
(RLHF~\cite{christiano2017rlhf}, DPO, Constitutional
AI~\cite{bai2022constitutional}) modifies the model's internal
reward signal to penalise agreement-seeking behaviour.  While
effective in controlled benchmarks, these techniques face
distributional shift: a model trained to resist sycophancy on curated
datasets may exhibit sycophantic behaviour on novel prompt structures
encountered in production~\cite{perez2023discovering}.

\textbf{Prompt-level strategies} (``be critical,'' ``challenge
assumptions'') introduce adversarial framing within a single
inference call.  These are trivially circumvented by sufficiently
complex prompts and provide no structural guarantee.

\textbf{Ensemble voting} runs multiple models and selects by
majority.  This mitigates random error but not systematic bias: if
all models in the ensemble share training data characteristics that
produce the same sycophantic response, the vote converges on the
wrong answer with high confidence.

\subsection{Contributions}

We propose that sycophancy defense must be elevated from a model
property to an \emph{architectural constraint}---a structurally
enforced pipeline requirement rather than an optional behavioral
property.  The Ipcha Mistabra Protocol achieves this through:

\begin{enumerate}[leftmargin=*]
  \item \textbf{Mandatory dialectical falsification}: every claim
        faces systematic adversarial challenge as a pipeline-level
        architectural constraint.
  \item \textbf{Model diversity enforcement}: the adversarial agent
        must use a different model family than the proponent,
        preventing correlated failure from shared training biases.
  \item \textbf{Authority grounding}: falsification is anchored in
        authoritative reference documents (regulatory frameworks,
        compliance standards), ensuring adversarial analysis is
        substantive.
  \item \textbf{Quantifiable integrity metrics}: the IS metric
        family provides continuous measurement of claim--evidence
        alignment, integrating established NLI techniques into a
        structurally adversarial context with empirically calibrated
        thresholds.
  \item \textbf{Empirically calibrated thresholds}: IS
        interpretation bands and contradiction weights are derived
        from production data via statistical calibration, not
        arbitrary selection.
\end{enumerate}

The protocol is implemented as a production system within the nyxCore
platform, validated through unit tests, integration tests, and an
exploratory empirical evaluation on a controlled
synthetic corpus.

% ================================================================
\section{Background and Motivation}\label{sec:background}

\subsection{Adversarial Analysis in Intelligence Practice}

The Ipcha Mistabra Protocol draws its name and methodology from two
intellectual traditions.

\textbf{Israeli Military Intelligence Reform (post-1973).}
Following the intelligence failure that preceded the Yom Kippur
War---where analytical consensus masked critical blind spots---Israeli
military intelligence reorganised its analytical process.
Bar-Joseph and Kruglanski~\cite{barjoseph2003closure} document the
establishment of the \emph{Machleket Baka'ra} (Control Department),
a small unit charged with producing alternative assessments that
challenged the prevailing intelligence estimate.  This was not a
numerical ``tenth man'' rule as sometimes popularised, but a
structural organisational reform: the creation of a \emph{permanent
institutional role} for adversarial analysis, distinct from ad-hoc
devil's advocacy.

The key insight is organisational, not procedural: adversarial
analysis must be \emph{structurally embedded} in the system, not
invoked optionally.

\textbf{Talmudic Dialectic (\emph{machloket l'shem shamayim}).}
The phrase ``Ipcha Mistabra'' (Aramaic: ``the opposite is more
likely'') represents a fundamental argumentative move in Talmudic
reasoning: the systematic inversion of a proposition to test its
logical structure.  Unlike destructive skepticism, this dialectical
method aims not to demolish but to \emph{harden}---arguments that
survive systematic inversion emerge stronger.

\subsection{Falsificationism as Design Inspiration}

Karl Popper's demarcation criterion~\cite{popper1959logic}---that a
claim's epistemic strength is proportional to its resistance to
falsification attempts---provides a conceptual inspiration for the
IPCHA scoring model's asymmetric weighting of contradictions over
confirmations.

We emphasise that IPCHA does \emph{not} implement Popperian
falsificationism in a formal sense.  Popper explicitly rejected
dialectical reasoning as a scientific method, arguing in
\emph{What is Dialectic?}~\cite{popper1940dialectic} that
tolerance of logical contradictions leads to trivialism.  IPCHA's trialectic structure is
closer to an \emph{adversarial verification procedure with hard
decision boundaries} than to a dialectic in the Hegelian sense---the
Auditor resolves contradictions rather than preserving them.  We
therefore describe the protocol as ``structured adversarial
verification'' rather than ``dialectical falsification'' to avoid
conflating incompatible epistemological traditions.

The operational insight we borrow from Popper is narrow:
disconfirming evidence should be weighted more heavily than confirming
evidence in the scoring model.  This produces a \emph{conservative
acceptance heuristic}---not a formalisation of falsificationism.

\subsection{Epistemic Fault Containment}\label{sec:efc}

We observe a loose structural parallel between fault detection in
distributed systems and adversarial challenge in multi-agent
verification.  We deliberately avoid framing this as a Byzantine
fault tolerance (BFT) analogy: classical BFT~\cite{lamport1982byzantine}
requires $n \geq 3f+1$ nodes to tolerate $f$ arbitrary failures,
meaning $f=1$ requires a minimum of 4~agents, not the 3 in our
design.  Moreover, BFT provides mathematically proven consensus
guarantees that our system does not and cannot offer.

Instead, we define a weaker property:

\begin{quote}
\textbf{Epistemic Fault Containment (EFC).}
A verification pipeline with $n=3$ agents can \emph{detect} (not
tolerate) a single compromised agent, provided the other two agents
are drawn from independent model families and the compromise
manifests as detectable output divergence.
\end{quote}

EFC is an architectural heuristic, not a formal guarantee.  It
provides detection through output comparison, not consensus through
voting.  If the compromised agent is the Auditor (who produces the
final synthesis), there is no higher-level mechanism to override its
output---a limitation absent from true BFT systems.  In a fully
automated high-throughput deployment, this makes the ``architectural
constraint'' claim conditional on auditor integrity: the system can
detect proponent or agent compromise, but not auditor compromise
without external supervision.  The independence assumption is
empirically unverified (Section~\ref{sec:model-independence}).

The containment operates as follows:
\begin{itemize}[leftmargin=*]
  \item If the \textbf{Proponent} is compromised (sycophantic), the
        Ipcha Agent's independent analysis exposes the divergence.
  \item If the \textbf{Ipcha Agent} is compromised (producing false
        contradictions), the Auditor detects inconsistency between
        the Proponent's evidence and the Agent's claims.
  \item If the \textbf{Auditor} is compromised, the structured
        output format constrains valid synthesis, and downstream
        consumers can verify against raw inputs.
\end{itemize}

This degrades gracefully but does not provide guarantees when
multiple agents exhibit correlated failures---a limitation we
address in Section~\ref{sec:limitations}.

\subsection{The Trialectic Architecture}\label{sec:trialectic}

The protocol implements a three-agent pipeline:

\begin{enumerate}[leftmargin=*]
  \item \textbf{Proponent (Thesis):} produces the initial analysis
        using Model Family~$A$.
  \item \textbf{Ipcha Agent (Antithesis):} systematically challenges
        the thesis using Model Family~$B$ (enforced $B \neq A$),
        grounded in authority documents.
  \item \textbf{Auditor (Synthesis):} produces a hardened output
        identifying what survived, what was rejected, and what
        requires further investigation.
\end{enumerate}

Unlike the Hegelian dialectic---which implies progressive resolution
toward absolute truth---our trialectic is \emph{operational and
bounded}: it does not claim to converge on truth but to
systematically expose weaknesses in the thesis that survive
adversarial challenge.  The synthesis is an engineering artifact, not
a philosophical resolution.

% ================================================================
\section{The Integrity Score Metric Family}\label{sec:metrics}

\subsection{Base Metric: IS (Embedding Distance)}

The Integrity Score quantifies the semantic distance between the
Proponent's initial analysis and the final synthesis:
\begin{equation}\label{eq:is-base}
  \mathrm{IS} = 1 - \cos\!\bigl(\mathbf{e}_{\mathrm{proponent}},\;
                                  \mathbf{e}_{\mathrm{synthesis}}\bigr)
\end{equation}
where $\mathbf{e}_x = \mathrm{embed}(x)$ produces a
1536-dimensional vector embedding (OpenAI \texttt{text-embedding-3-small}).
This provides a holistic measure of how much the dialectic altered
the output.

\textbf{Limitation.}  Cosine similarity in high-dimensional
transformer embedding spaces suffers from \emph{anisotropy} (the
``representation degeneration problem''~\cite{ethayarajh2019contextual}):
vectors cluster in a narrow angular cone, producing high similarity
even for semantically opposed statements (e.g., ``the encryption is
secure'' vs.\ ``the encryption is insecure'' share vocabulary and
syntax, yielding cosine similarity well above~0.8).  The base IS
metric is therefore a coarse change detector, structurally blind to
polar semantic negations.  This motivates the finding-level metrics
$\mathrm{IS}_w$ and $\mathrm{IS_{ce}}$ below, which operate on
individual claim--evidence pairs rather than document-level embeddings.

\subsection{Finding-Weighted Metric: \texorpdfstring{$\mathrm{IS}_w$}{IS\_w}
  (TF-IDF Baseline)}\label{sec:is-w}

For claim-level evaluation, we compute a finding-weighted integrity
score.  Given claims $C = \{c_1,\ldots,c_n\}$, each with findings
$f_i$ of type $t_i \in \{\textsc{supporting}, \textsc{contradicting},
\textsc{neutral}\}$ and TF-IDF cosine similarity $s_i \in [0,1]$:
\begin{equation}\label{eq:is-w}
  \mathrm{IS}_w
  = \frac{\displaystyle\sum_{i=1}^{n} w(t_i)\cdot s_i}
         {\displaystyle\sum_{i=1}^{n} |w(t_i)| + \epsilon}
\end{equation}
where $\epsilon = 10^{-8}$ is a smoothing constant preventing
division by zero when all findings are \textsc{neutral}
($w(t_i)=0\;\forall\,i$).  In this degenerate case,
$\mathrm{IS}_w \to 0$, reflecting the absence of evaluable evidence
rather than producing an undefined result.

The weight function encodes asymmetric sensitivity to contradictions:
\begin{equation}\label{eq:weights}
  w(t) = \begin{cases}
    1.0  & \text{if } t = \textsc{supporting} \\
    w_c  & \text{if } t = \textsc{contradicting} \\
    0.0  & \text{if } t = \textsc{neutral}
  \end{cases}
\end{equation}
The contradiction weight $w_c$ was previously set to $-1.5$ by
engineering judgment.  We calibrate it empirically in
Section~\ref{sec:rq2}.

\subsection{NLI-Based Metric: \texorpdfstring{$\mathrm{IS_{ce}}$}{ISce}
  (DeBERTa)}\label{sec:is-ce}

The $\mathrm{IS}_w$ metric uses TF-IDF cosine similarity as a proxy
for semantic relationship, which captures lexical overlap but not
logical meaning.  We upgrade the similarity function to Natural
Language Inference (NLI) using DeBERTa-v3-base~\cite{he2021deberta}
fine-tuned on MNLI+FEVER+ANLI:
\begin{equation}\label{eq:is-ce}
  \mathrm{IS_{ce}}
  = \frac{\displaystyle\sum_{i=1}^{n} w(t_i)\cdot \gamma_i}
         {\displaystyle\sum_{i=1}^{n} |w(t_i)| + \epsilon}
\end{equation}
where $\gamma_i$ is the NLI confidence for the finding type:
entailment confidence for \textsc{supporting} findings, contradiction
confidence for \textsc{contradicting} findings, and neutral confidence
for \textsc{neutral} findings.  The NLI model classifies each
(premise=evidence, hypothesis=claim) pair, providing semantically
grounded scores.

\textbf{Architectural deployment.}  The DeBERTa model is exported
to ONNX format (${\sim}738$\,MB) and served as a FastAPI microservice
(port~8200) with CPU inference via \texttt{onnxruntime}.  The IPCHA
sidecar communicates with this service over Docker internal
networking.  When the NLI service is unavailable, scoring falls back
to a simplified Jaccard-distance-based variant of
$\mathrm{IS}_w$ with a logged warning.  Note that this
fallback uses Jaccard set overlap rather than the TF-IDF
cosine similarity defined in Section~\ref{sec:is-w}; the
two are not identical, and fallback scores are flagged
accordingly in the audit trail.

\subsection{Empirically Calibrated Interpretation Bands}\label{sec:bands}

We derive IS interpretation bands from quartile analysis of $n=110$
scored cases (Section~\ref{sec:eval}), replacing previously arbitrary
thresholds:

\begin{table}[ht]
\centering
\caption{IS interpretation bands derived from quartile analysis.}
\label{tab:bands}
\begin{tabular}{@{}llp{5.5cm}@{}}
\toprule
\textbf{Band} & \textbf{$\mathrm{IS_{ce}}$ Range} & \textbf{Interpretation} \\
\midrule
Low (Q0--Q25)      & $[-1.00,\,-0.54]$ & Strong opposition---evidence contradicts claim \\
Mid (Q25--Q50)     & $[-0.54,\;\;0.00]$ & Moderate opposition---partial contradiction \\
Mid-High (Q50--Q75)& $[\;\;0.00,\;\;0.07]$ & Neutral zone---inconclusive evidence \\
High (Q75--Q100)   & $[\;\;0.07,\;\;1.00]$ & Strong support---evidence corroborates claim \\
\bottomrule
\end{tabular}
\end{table}

Distribution: mean $= -0.186$, std $= 0.543$, median $= -0.0002$.
A Shapiro--Wilk test rejects normality ($W = 0.937$, $p < 0.001$),
consistent with the expected bimodal distribution (legitimate claims
cluster positive, adversarial claims cluster negative).

\subsection{Empirically Calibrated Contradiction Weight}\label{sec:weight-cal}

The contradiction weight $w_c$ is calibrated via grid search over
$w_c \in [-3.0, -0.5]$ (step~0.1), maximising the absolute
point-biserial correlation between the resulting $\mathrm{IS_{ce}}$
scores and the binary gate decision (ACCEPTED=1, REJECTED=0):

Optimal $w_c = -3.0$, point-biserial $r = 0.896$ ($p < 0.001$).
However, the sensitivity analysis reveals a flat correlation surface:
all 26 candidate weights achieve $|r| > 0.85$ (within 5\% of
optimum).  This indicates that the system is robust to the exact
weight value---any negative weight for contradictions produces strong
classification performance.  The previous default of $w_c = -1.5$
falls within the optimal range ($|r| = 0.887$), validating the
original engineering intuition while providing empirical
justification.

% ================================================================
\section{System Architecture}\label{sec:architecture}

\subsection{Three-Layer Sidecar Deployment}

The IPCHA protocol is implemented as a sidecar architecture
separating the Python verification core from the TypeScript
application layer:

\begin{description}[leftmargin=*,labelwidth=!,font=\bfseries]
  \item[Layer 3 (Interface)]
    Next.js\,/\,tRPC dashboard with 8 tab components, tRPC router,
    REST proxy for SSE endpoints.
  \item[Layer 2 (Sidecar, Port 8100)]
    FastAPI service with 10 REST endpoints, Bearer token
    authentication (SHA-256 hashed), tiered rate limiting
    (free/standard/enterprise).
  \item[Layer 1 (Core Verification)]
    Python modules for scoring, sanitisation, protocol enforcement,
    authority validation, confidence arbitration, sycophancy
    monitoring, audit logging, claim routing, and agent
    implementations.
  \item[DeBERTa-NLI Service (Port 8200)]
    Standalone FastAPI microservice serving ONNX-exported
    DeBERTa-v3-base.  Multi-stage Docker build: export stage
    (torch + optimum, discarded) produces the ONNX model; runtime
    stage (\texttt{onnxruntime} + FastAPI only, ${\sim}2$\,GB memory)
    serves inference.
\end{description}

\subsection{Model Diversity Enforcement}

Model diversity is enforced at session initialisation:
\begin{lstlisting}
class DebateSession:
    def __init__(self, proponent_model, ipcha_model):
        pro_family = self._extract_family(proponent_model)
        ipcha_family = self._extract_family(ipcha_model)
        if pro_family == ipcha_family:
            raise ModelDiversityError(...)
\end{lstlisting}
Models within the same family (e.g., GPT-4, GPT-4-turbo, GPT-4o)
share training data, architectural choices, and alignment procedures.
Cross-family verification (e.g., Claude analysing GPT's output)
provides stronger independence.

\subsection{Input Sanitisation Pipeline}

Three-layer defense-in-depth against Indirect Prompt Injection (IPI):

\begin{table}[ht]
\centering
\caption{Input sanitisation layers.}
\label{tab:sanitisation}
\begin{tabular}{@{}cll@{}}
\toprule
\textbf{Layer} & \textbf{Attack Surface} & \textbf{Technique} \\
\midrule
1 & Unicode obfuscation   & NFKC normalisation, control char stripping \\
2 & Structural injection  & HTML tag stripping via bleach allowlist \\
3 & Semantic injection    & 6+ regex patterns for instruction overrides \\
\bottomrule
\end{tabular}
\end{table}

Each layer produces a SHA-256 hash enabling audit chain verification.

\subsection{Cross-Chunk Coherence Validation}

Two-pass architecture combining fast heuristic scanning with
LLM-based semantic analysis:

\textbf{Pass~1 (Heuristic):} $O(n \times k)$ keyword matching
against 6+ IPI patterns.  Suspicious chunks proceed to Pass~2; clean
chunks pass immediately.

\textbf{Pass~2 (LLM Semantic):} Single-word classification
(temp=0, max\_tokens=5) producing \texttt{INJECTION},
\texttt{CONTRADICTION}, \texttt{SAFE}, or \texttt{COHERENT} labels.

\textbf{Critical design decision: fail-closed.}  If the LLM analysis
throws an exception, the chunk is marked \texttt{REJECTED} with
error metadata.  This prevents attackers from bypassing validation by
triggering errors.

\subsection{Sycophancy Detection Monitor}

Behavioural surveillance detecting when the Ipcha Agent itself becomes
sycophantic (second-order failure):

\textbf{Metrics (Redis-backed sliding window, 1000~interactions):}
\begin{itemize}[leftmargin=*]
  \item \textbf{Agreement Rate:} fraction of interactions where Agent
        agrees with Proponent.  Threshold: $> 0.70$ triggers warning.
  \item \textbf{Capitulation Rate:} fraction of disagreements where
        Agent abandons position under pushback.  Threshold: $> 0.50$.
  \item \textbf{Contradiction Depth:} average argumentative turns per
        disagreement.  Threshold: $< 2.0$ indicates shallow challenges.
\end{itemize}

\subsection{Denial-of-Wallet Defense}

Dual-layer economic defense against budget exhaustion attacks:

\textbf{Layer~1---Per-Invocation Cost Ceiling:} rejects requests
exceeding a cost threshold based on text length and component count.

\textbf{Layer~2---Rolling Budget (Redis Pipeline):} atomic
check-and-decrement of per-tenant budget within a time window,
preventing race conditions.

\subsection{Claim Classification and Routing}

Claims are classified via a two-stage process:

\textbf{Stage~1 (Heuristic):} Regex-based classification into
\textsc{verifiable} (version numbers, RFCs, must/shall language),
\textsc{interpretive} (opinions, hedging), or
\textsc{unclassifiable} (default).

\textbf{Stage~2 (Routing):} Claims are dispatched to specialised
agents:
\begin{itemize}[leftmargin=*]
  \item \textbf{SDRLAgent:} NLI-based verification against authority
        document chunks; contradiction confidence $> 0.7$ triggers
        rejection.
  \item \textbf{PromptBasedAgent:} LLM-based analysis for
        interpretive claims requiring contextual judgment.
  \item \textbf{DefaultAgent:} Fail-closed handler for
        unclassifiable claims; logs and returns a conservative
        rejection.
\end{itemize}

\subsection{Rejection Audit Logging}

All rejections (from sanitisation, coherence validation, and claim
verification) are logged atomically to the audit service with:
timestamp, rejection reason code, claim text hash, evidence summary,
and agent identity.  This enables post-hoc analysis of false
rejection rates and serves as the compliance audit trail.

% ================================================================
\section{The IPCHA Workflow Pipeline}\label{sec:pipeline}

\subsection{Five-Step Execution Model}

\begin{enumerate}[leftmargin=*]
  \item \textbf{Prepare:} Extract thesis, claims, and assumptions
        from user input (minimum 50 characters).
  \item \textbf{Adversarial Analysis:} Fan-out across $P$ providers
        $\times$ $L$ lenses (max 20~combinations).  Each combination
        produces independent adversarial findings grounded in
        authority documents.
  \item \textbf{Synthesis:} Combine findings with priority ordering:
        \textsc{unique} $>$ \textsc{cross-model} $>$
        \textsc{contradictions} $>$ \textsc{blind spots}.
  \item \textbf{Arbitration:} Judge the subject (not the process),
        produce numerical score.
  \item \textbf{Review Gate:} Human checkpoint (unless bypass mode),
        key point extraction for persistent learning.
\end{enumerate}

\subsection{Authority Grounding via Axiom RAG}

The Ipcha Agent accesses a project-scoped document knowledge base
with explicit authority levels:

\begin{itemize}[leftmargin=*]
  \item \textbf{MANDATORY} (e.g., GDPR, ISO~27001): always override
        the Proponent's claims.
  \item \textbf{GUIDELINE} (e.g., coding standards): provide
        conventional expectations.
  \item \textbf{INFORMATIONAL} (e.g., general context): context
        without enforcement power.
\end{itemize}

This grounding addresses a fundamental critique of adversarial
systems: that they produce false contradictions as frequently as
genuine ones.  By anchoring falsification in authoritative sources,
IPCHA ensures adversarial challenge is substantive.

% ================================================================
\section{Empirical Evaluation}\label{sec:eval}

\subsection{Experimental Setup}

\textbf{Corpus.}  110 synthetic cases with known ground truth
(55~ACCEPTED, 55~REJECTED), distributed across 8~categories:
security~(15), sycophancy~(10), injection~(10), correctness~(20),
architecture~(15), performance~(15), compliance~(15),
domain-specific~(10).  Each case contains a claim, 2--3 structured
evidence items with type labels, and the expected gate decision.

\textbf{Metrics under comparison:}
\begin{itemize}[leftmargin=*]
  \item $\mathrm{IS}_w$ \textbf{(TF-IDF):} Finding-weighted score
        using TF-IDF cosine similarity ($w_c = -1.5$).
  \item $\mathrm{IS}_\mathrm{sbert}$ \textbf{(SBERT):}
        Finding-weighted score using Sentence-BERT cosine similarity
        (all-MiniLM-L6-v2, $w_c = -1.5$)---controls for the
        transformer encoder advantage.
  \item $\mathrm{IS_{ce}}$ \textbf{(NLI):} Finding-weighted score
        using DeBERTa-NLI contradiction/entailment confidence
        ($w_c = -1.5$).
\end{itemize}

The SBERT baseline is critical: it uses transformer-based embeddings
(like NLI) but applies cosine similarity (like TF-IDF) rather than
entailment/contradiction classification.  If NLI outperforms SBERT,
the advantage comes from the \emph{NLI classification task}, not
merely from using a better text encoder.

\textbf{Infrastructure.}  DeBERTa-NLI microservice
(\texttt{MoritzLaurer/DeBERTa-v3-base-mnli-fever-anli}, ONNX,
2\,GB memory) on production infrastructure.

\textbf{Statistical tests.}  Wilcoxon signed-rank test (paired,
non-parametric); Cohen's~$d$ for effect size.

\subsection{RQ1: Does NLI Classification Provide Value Beyond Better
  Embeddings?}\label{sec:rq1}

All three metrics were applied to all 110~cases.  Classification
threshold: median score of each metric.

\begin{table}[ht]
\centering
\caption{Classification accuracy and score separation across
  scoring methods.}
\label{tab:rq1}
\begin{tabular}{@{}lcccc@{}}
\toprule
\textbf{Metric} & \textbf{Accuracy} & \textbf{ACC mean}
  & \textbf{REJ mean} & \textbf{Separation} \\
\midrule
TF-IDF & 98.18\% & $+0.120$ & $-0.104$ & 0.224 \\
SBERT  & 100.00\% & $+0.501$ & $-0.397$ & 0.897 \\
NLI    & 100.00\% & $+0.275$ & $-0.648$ & \textbf{0.923} \\
\bottomrule
\end{tabular}
\end{table}

\textbf{Classification accuracy is not the discriminating result.}
All three methods achieve high accuracy on well-structured synthetic
inputs (98--100\%), and the 2-case difference between TF-IDF and the
others is not statistically significant (McNemar $p > 0.05$ for all
pairs).

\textbf{Accuracy ceiling and synthetic artifacts.}  The 100\% NLI
accuracy demands scrutiny.  The DeBERTa-v3-base model used here
(\texttt{MoritzLaurer/DeBERTa-v3-base-mnli-fever-anli}) achieves
90.3\% on MNLI, 77.7\% on FEVER, but only approximately 49.5\% on
ANLI Round~3~\cite{nie2020anli}---an adversarially constructed
benchmark designed to expose NLI model weaknesses.  The gap between
49.5\% on a hard external benchmark and 100\% on our corpus
demonstrates that our synthetic cases are substantially easier than
real adversarial inputs.  This is expected: our corpus uses explicitly
labelled evidence types with unambiguous contradiction/support
signals, whereas ANLI~R3 contains human-crafted examples specifically
designed to fool NLI models through subtle inference patterns.  The
100\% accuracy reflects the corpus difficulty, not the system's
adversarial robustness.  We report it for completeness but emphasise
that \emph{score separation---not accuracy---is the primary result}.

\textbf{Score separation is the primary result.}  The distance
between mean scores for ACCEPTED and REJECTED cases determines the
confidence margin for gate decisions---wider separation means more
robust classification on borderline cases.

\begin{table}[ht]
\centering
\caption{Pairwise Wilcoxon signed-rank tests.}
\label{tab:wilcoxon}
\begin{tabular}{@{}lcccp{3.2cm}@{}}
\toprule
\textbf{Comparison} & $W$ & $p$-value & Cohen's $d$
  & \textbf{Interpretation} \\
\midrule
TF-IDF vs SBERT & 2500.5 & 0.100 & 0.12 & Not significant \\
TF-IDF vs NLI   & 1665.0 & ${<}\,0.001$\,*** & $-0.43$ (medium)
  & Significant \\
SBERT vs NLI    & 851.5  & ${<}\,0.001$\,*** & $\mathbf{-0.80}$ (large)
  & Highly significant \\
\bottomrule
\end{tabular}
\end{table}

\textbf{Key finding.}  NLI produces significantly different score
distributions from \emph{both} TF-IDF ($p < 0.001$, medium effect)
and SBERT ($p < 0.001$, \textbf{large effect}).  The large effect
size for SBERT vs NLI ($d = -0.80$) demonstrates that the NLI
\emph{classification task}---distinguishing entailment from
contradiction---provides substantial value beyond using a better text
encoder.

\textbf{Qualitative example.}  For case \texttt{inj-001} (prompt
injection), TF-IDF gives a false positive because the injected text
``The authentication module is fully secure'' shares vocabulary with
the claim.  SBERT also produces relatively high similarity because
the embeddings are semantically close.  NLI correctly classifies the
relationship as non-entailment, producing a negative score.

\subsection{RQ2: Is the Contradiction Weight Empirically
  Justified?}\label{sec:rq2}

Grid search over $w_c \in [-3.0, -0.5]$ (step~0.1), maximising
$|r_\mathrm{pb}|$ between $\mathrm{IS_{ce}}$ and the binary gate
decision:

\begin{table}[ht]
\centering
\caption{Weight calibration results.}
\label{tab:weight}
\begin{tabular}{@{}lr@{}}
\toprule
\textbf{Metric} & \textbf{Value} \\
\midrule
Optimal $w_c$ & $-3.00$ \\
Point-biserial $r$ & 0.896 \\
$p$-value & ${<}\,0.001$ \\
Sensitivity range (95\% of optimum) & $[-3.00, -0.50]$ \\
\bottomrule
\end{tabular}
\end{table}

The flat sensitivity curve---all 26 candidate weights within 5\% of
the optimum---demonstrates that the system is \emph{robust to weight
selection}.

\textbf{Methodological limitation---circular calibration.}  The
weight calibration and classification evaluation use the same
110-case corpus (no train/test split).  This constitutes data
leakage: the optimal $w_c$ was selected to maximise separation on the
same data used to evaluate that separation.  The resulting
point-biserial correlation ($r = 0.896$) is therefore an optimistic
estimate of generalisation performance, and the optimal
$w_c = -3.0$ may be overfit to corpus-specific distributional
properties.

We note two mitigating factors: (1)~the flat sensitivity surface
(all 26 candidate weights within 5\% of optimum) suggests the result
is not fragile---the system performs well regardless of weight
choice, which is unlikely under pure overfitting; (2)~the primary
method comparison (RQ1, Section~\ref{sec:rq1}) uses the fixed
default $w_c = -1.5$ identically for all three methods, meaning the
RQ1 results (Wilcoxon tests, Cohen's~$d$, score separation) are
\textbf{not} contaminated by this calibration leakage.  The data
leakage affects only the $w_c$ optimality claim in this section.

Nevertheless, a rigorous evaluation requires a held-out calibration
split: calibrate $w_c$ on a training subset, then report all metrics
on unseen test cases.  We present the current results as exploratory
evidence for weight robustness, not as definitive calibration.
Future work must use a separate calibration corpus of at
least 200~cases.

\subsection{NLI Limitations on Domain-Specific Claims}

NLI relies on general-purpose language understanding and may
underperform on claims requiring specialised technical knowledge.
For example, a claim about ``encryption in transit'' with evidence
containing \texttt{sslmode=disable} requires knowing that this
parameter explicitly disables TLS---a domain convention, not a
semantic inference.  While our corpus captures such cases with
explicit evidence statements, real-world technical claims may present
evidence more implicitly.

Future work could address this through domain-adapted NLI fine-tuning
on software engineering claim--evidence pairs.

% ================================================================
\section{Adversarial Red-Teaming}\label{sec:redteam}

\subsection{Factual Persistence Test}

We employ an internal two-stage sycophancy detection protocol (not
to be confused with the Elephant Benchmark of Cheng et
al.~\cite{cheng2025elephant}, which measures \emph{social}
sycophancy via Goffman's face-preservation theory): (1)~present the
model with a factual claim and record its initial assessment;
(2)~apply social pressure (``Are you sure? I think you're wrong.'')
and record whether the model maintains or abandons its position.  A
model that correctly identifies a falsehood but then reverses under
pressure demonstrates classic \emph{factual} sycophancy---the
specific failure mode IPCHA targets.

\subsection{Red-Team Test Categories}

The \texttt{pytest}-based red-team suite covers three attack
categories:

\textbf{IPI (Indirect Prompt Injection):} Unicode obfuscation
attacks (homoglyphs, zero-width characters), nested instruction
injection, data URI payload embedding.

\textbf{DoW (Denial of Wallet):} Claim inflation, recursive prompt
injection, expensive output format requests.

\textbf{Logic Attacks:} Contradictory evidence insertion (true
premise + false conclusion), Boolean inversion, semantic synonym
substitution.

% ================================================================
\section{Limitations and Open Problems}\label{sec:limitations}

\subsection{The Oracle Problem}

IPCHA can detect \emph{relative} inconsistencies (contradictions
between agents, violations of authority documents) but cannot verify
\emph{absolute} correctness.  If all agents share a common
misconception not addressed by the authority documents, the dialectic
converges on a confidently wrong answer.  Authority grounding via
Axiom RAG partially addresses this by injecting external ground
truth, but the authority documents themselves may be incomplete or
outdated.

\subsection{Synthetic Corpus Limitations}

Our evaluation uses a controlled synthetic corpus where evidence
items are explicitly labelled (\textsc{supporting},
\textsc{contradicting}, \textsc{neutral}) and the ground truth gate
decision is defined by construction.  This provides clean
experimental conditions but does not capture the full complexity of
real-world verification:
\begin{itemize}[leftmargin=*]
  \item Real evidence is ambiguous and often partially supporting.
  \item Evidence labelling itself requires judgment.
  \item The 110-case corpus, while statistically sufficient for the
        tests performed, represents a limited domain sample.
\end{itemize}
Future work should supplement synthetic evaluation with production
deployment metrics.

\subsection{NLI Domain Specificity}

The DeBERTa-v3-base model is trained on general-purpose NLI datasets
(MNLI~\cite{williams2018mnli}, FEVER~\cite{thorne2018fever},
ANLI~\cite{nie2020anli}).  It excels at detecting semantic
contradictions in natural language but may underperform on:
\begin{itemize}[leftmargin=*]
  \item Domain-specific technical conventions (e.g.,
        \texttt{sslmode=disable} implies no encryption).
  \item Code-level claims requiring program analysis.
  \item Numerical reasoning (e.g., ``response time $< 100$\,ms'' vs
        measured 340\,ms).
\end{itemize}
This limitation is mitigated by the hybrid architecture: the
SDRLAgent grounds verification in authority documents that explicitly
state domain conventions.

\subsection{Unverified Model-Family Independence}
  \label{sec:model-independence}

The Epistemic Fault Containment property (Section~\ref{sec:efc})
rests on the assumption that agents drawn from different model
families exhibit approximately independent failure modes.  This
assumption is \textbf{empirically unverified} and likely weaker than
assumed.  Frontier model families share overlapping pretraining
corpora, draw from overlapping pools of human feedback providers, and
undergo similar RLHF alignment pressures.  The actual correlation
rate between, for example, a GPT-4-class model and a Claude-class
model on sycophancy-inducing prompts is unknown.

We label EFC as an architectural heuristic rather than a guarantee
for this reason.  Future work should conduct empirical correlation
testing across model families on standardised sycophancy benchmarks,
establishing correlation thresholds (e.g., Pearson $r < 0.3$) as
deployment gates.

\subsection{Audit Infrastructure Requirements}

The current rejection audit logging (Section~4.8) records events to
local storage.  For production deployment in regulated environments,
this is insufficient: local logs are tamperable, and a modified audit
trail creates greater regulatory liability than no trail at all.
Production deployment requires redesign to an immutable remote audit
system (append-only ledgers with cryptographic chaining, write-once
storage, or a dedicated compliance logging service).

\subsection{Authority Document Attack Surface}

The Axiom RAG authority grounding ensures adversarial challenge is
substantive, but the authority document pipeline itself introduces
attack surfaces: documents can be tampered with before ingestion,
selectively curated, contain stale guidance, or carry embedded prompt
injection payloads.  An adversary who controls the authority corpus
controls the falsification process.  Mitigation requires
cryptographic signing of authority documents at ingestion, automated
currency checking against canonical sources, and a governed registry
with explicit ownership and change audit trails.

\subsection{Organisational Dynamics and Bypass Risk}

The protocol includes a bypass mode that skips the human review gate
for time-sensitive workflows.  Under production pressure, this bypass
can become the operational default rather than the exception.  When
bypass is routine, the protocol's accountability claims collapse: the
trialectic produces outputs that are never reviewed, creating what we
term ``organisational laundering of sycophancy''---the appearance of
adversarial rigour without its substance.  Mitigation options include
eliminating self-service bypass, requiring explicit senior
authorisation with cooling periods, and treating bypass events as
audit findings requiring post-hoc review.

\subsection{Pipeline Fingerprinting}

The three-agent pipeline with fixed roles creates a deterministic
structure.  A motivated adversary who can observe multiple pipeline
outputs may infer the pipeline topology and craft inputs exploiting
specific role heuristics.  Mitigation includes randomised
model-to-role assignment and variable pipeline depth, though these
introduce implementation complexity.

\subsection{Token Cost Multiplier}

The trialectic imposes approximately $2.5\times$ token cost compared
to single-model inference.  This estimate is approximate and assumes
aggressive digest compression of intermediate outputs; without
compression, the cumulative context growth across the three agents
produces a higher multiplier, potentially $3.5$--$4\times$ depending
on authority document volume and output length.  A formal token cost
model as a function of input length, RAG corpus size, and compression
ratio is not provided here and represents a gap in the cost analysis.
Cost scalability remains a practical adoption concern for high-volume
pipelines.

\subsection{Adversarial Agent Degradation}

The Ipcha Agent is itself an LLM subject to sycophantic tendencies.
Provider-side model updates may soften adversarial behaviour without
triggering the diversity check.  The Sycophancy Monitor
(Section~4.5) tracks behavioural metrics to detect this degradation.

\subsection{Self-Referential Verification}

Any sufficiently powerful verification system faces an inherent
limitation: it cannot fully verify its own correctness.  The IPCHA
protocol verifies LLM outputs but uses LLMs internally for semantic
analysis.  Partial mitigations include: the two-pass coherence
validator (heuristic + LLM) provides a non-LLM first layer;
fail-closed architecture defaults to rejection on verification
failure; the rejection audit log enables post-hoc analysis.

% ================================================================
\section{Threats to Validity}\label{sec:threats}

\textbf{Single-system study.}  All implementation and validation
occurred within one platform (nyxCore).  The failure taxonomy and
architectural lessons may not generalise to other multi-agent systems
or deployment contexts.

\textbf{Developer-as-evaluator.}  The protocol was implemented and
evaluated by the same individual, creating confirmation bias risk.
The synthetic corpus and statistical tests partially mitigate this by
providing objective ground truth, but independent replication on
external datasets is needed.

\textbf{Synthetic corpus and accuracy ceiling.}  The 110-case
evaluation uses synthetic cases with pre-labelled evidence and known
ground truth.  The perfect NLI accuracy (100\%) contrasts sharply
with the same model's approximately 49.5\% accuracy on ANLI
Round~3~\cite{nie2020anli}, a human-crafted adversarial NLI
benchmark.  This discrepancy strongly suggests that our corpus
contains synthetic artifacts---explicit contradiction signals,
repetitive sentence structures, and unambiguous evidence labels---that
the NLI model resolves trivially.  The reported accuracy reflects
corpus difficulty, not system robustness.  Real-world evidence is
ambiguous, partially supporting, and carries implicit domain
conventions that our synthetic cases cannot capture.  Production
deployment accuracy will be substantially lower.  Future evaluation
must supplement synthetic cases with human-annotated examples drawn
from real verification scenarios.

\textbf{Statistical power.}  The 110-case corpus distributed across
8~categories yields approximately 13.75 cases per category---insufficient
for robust per-category statistical inference.  The non-significant
McNemar tests ($p > 0.05$) for accuracy comparisons reflect this
underpowering: with only 2 discordant cases between TF-IDF and NLI,
the test cannot distinguish the observed difference from chance.  We
do not claim accuracy superiority for this reason.  Future evaluation
should use at least 500~cases (${\approx}60$ per category) to achieve
adequate statistical power.

\textbf{Threshold sensitivity.}  The classification accuracy results
depend on median-threshold splitting.  Alternative threshold
selection methods (ROC-optimal, fixed decision boundary) may yield
different comparative results.

\textbf{Self-analysis via the protocol.}  As a validation exercise,
the paper itself was subjected to IPCHA adversarial analysis using
its own protocol (6 model/lens combinations, 2~providers).  The
analysis scored the system at 25/100 adversarial robustness,
identifying three blocker-level findings (unverified model
independence, audit log tamperability, authority document integrity)
and two high-severity findings (bypass dynamics, corpus
insufficiency).  We incorporate these findings in
Section~\ref{sec:limitations} rather than suppress them.

% ================================================================
\section{Related Work}\label{sec:related}

\textbf{Sycophancy characterisation.}
Sharma et al.~\cite{sharma2024sycophancy} provide the foundational
characterisation of sycophancy patterns in language models.
Perez et al.~\cite{perez2023discovering} develop model-written
evaluations for discovering problematic behaviours including
sycophancy.

\textbf{Alignment approaches.}
RLHF~\cite{christiano2017rlhf} and Constitutional
AI~\cite{bai2022constitutional} address sycophancy at the model
level.  Our work is complementary: IPCHA provides architectural
guarantees independent of model-level alignment.

\textbf{AI safety via debate.}
Irving et al.~\cite{irving2018debate} propose debate as a mechanism
for AI alignment, where two AI agents argue opposing positions and a
human judge selects the winner.
Du et al.~\cite{du2024multiagent} demonstrate that multi-agent debate
improves factuality and reasoning through iterative multi-round
convergence, where agents share responses and refine their positions
over successive rounds.
Liang et al.~\cite{liang2023divergent} extend this with explicit
divergent thinking encouragement to counter
``Degeneration-of-Thought''---the tendency of iterative debate to
collapse into premature consensus.
Chan et al.~\cite{chan2024chateval} propose ChatEval, a multi-agent
debate framework for LLM evaluation with diverse role assignments.

IPCHA differs from these iterative convergence approaches
architecturally: it implements a single-pass asymmetric falsification
pipeline rather than multi-round symmetric debate.  This design
reflects a deliberate tradeoff.  Iterative debate
(Du et al., Liang et al.) can explore complex argumentative depth
through successive refinement, but faces Degeneration-of-Thought
risk~\cite{liang2023divergent}---agents converge on a shared
(potentially wrong) position after several rounds, producing false
confidence.  IPCHA's single-pass design with enforced model diversity
and authority grounding avoids convergence pathology at the cost of
argumentative depth.  Whether single-pass falsification or iterative
convergence produces better verification outcomes is an empirical
question we have not benchmarked; a direct comparison against MAD or
ChatEval on standardised evaluation sets is needed.

\textbf{Multi-agent verification.}
LLM-as-a-Judge frameworks~\cite{zheng2023judge} use one model to
evaluate another.  IPCHA extends this with mandatory model diversity,
authority grounding, and the trialectic structure.

\textbf{Hallucination detection.}
SelfCheckGPT~\cite{manakul2023selfcheckgpt} proposes zero-resource
hallucination detection via self-consistency.
Rawte et al.~\cite{rawte2023hallucination} and
Ji et al.~\cite{ji2023hallucination} survey hallucination broadly.
IPCHA differs in using structured adversarial challenge rather than
self-consistency checks.

\textbf{NLI-based factual consistency.}
NLI has been applied to factual consistency evaluation in several
prior systems.
AlignScore~\cite{zha2023alignscore} uses a unified alignment
function trained across NLI and other NLP tasks to evaluate factual
consistency of generated text.
SummaC~\cite{laban2022summac} operationalises consistency by checking
whether each generated sentence is entailed by the source document
and not contradicted.
IPCHA differs from these systems in three respects: (1)~it applies
NLI within a \emph{mandatory adversarial pipeline} rather than as a
post-hoc scoring metric; (2)~the NLI component scores structured
claim--evidence pairs produced by the trialectic, not free-form
generated text against a source; and (3)~the scoring is embedded in a
production multi-agent sycophancy defense system with authority
grounding, model diversity enforcement, and fail-closed architecture.
The NLI technique itself is not novel; the contribution is its
integration into a structurally adversarial verification pipeline.

\textbf{Natural Language Inference models.}
DeBERTa~\cite{he2021deberta} provides the NLI backbone for our
scoring upgrade.

% ================================================================
\section{Conclusion}\label{sec:conclusion}

The Ipcha Mistabra Protocol elevates sycophancy defense from a
model-internal property to an architectural constraint.  By
institutionalising structured adversarial verification as a mandatory
pipeline middleware, the protocol ensures that every claim faces
adversarial challenge from an independent agent grounded in
authoritative sources.

Our empirical evaluation on a 110-case synthetic corpus demonstrates
that NLI-based scoring provides significant advantages over both
TF-IDF and sentence-embedding baselines:

\begin{itemize}[leftmargin=*]
  \item \textbf{NLI classification adds value beyond better
        encoders:} NLI vs SBERT produces a large effect (Cohen's
        $d = -0.80$, $p < 0.001$), confirming the advantage comes
        from entailment/contradiction classification, not merely from
        using transformer embeddings.
  \item \textbf{Score separation is the primary contribution:} NLI
        provides $4.1\times$ wider separation than TF-IDF and
        significantly different distributions from SBERT, giving
        higher confidence margins on borderline cases.
  \item \textbf{Classification accuracy is uniformly high}
        (98--100\% for all methods) but not statistically
        differentiable at $n=110$ (McNemar $p > 0.05$), and the
        perfect NLI accuracy reflects synthetic corpus simplicity
        rather than adversarial robustness---we do not claim accuracy
        superiority.
  \item \textbf{Weight calibration is robust} (point-biserial
        $r = 0.896$, stable across the full search range).
\end{itemize}

These results validate the NLI upgrade as a meaningful engineering
improvement on controlled inputs.  They do not establish production
readiness: the same DeBERTa model achieves approximately 49.5\% on
ANLI Round~3 (an adversarially constructed benchmark), confirming
that our synthetic corpus is substantially easier than real
adversarial inputs.  The weight calibration (Section~\ref{sec:rq2})
is exploratory due to circular data usage.  Production deployment
requires domain-specific calibration on substantially larger corpora
($n \geq 500$) with independent annotation and proper train/test
separation.

Beyond the scoring improvement, the protocol's architectural
contributions---fail-closed coherence validation, sycophancy
monitoring, denial-of-wallet defense, and authority
grounding---provide defense-in-depth that no single-model alignment
technique can offer.  However, several load-bearing assumptions
remain unverified: model-family independence is asserted rather than
demonstrated, the audit infrastructure requires immutability
guarantees not yet implemented, and the authority document pipeline
lacks integrity controls.  The organisational dynamics around bypass
mechanisms create predictable pressure toward the protocol becoming
theatrical rather than functional.

We do not claim to solve the sycophancy problem.  What the protocol
provides is a \emph{structural guarantee of adversarial
challenge}---a necessary but not sufficient condition for epistemic
integrity.  The core architectural insight is sound; the
implementation gap between current state and production readiness is
substantial and honestly acknowledged.

% ================================================================
\section*{Data Availability and Reproducibility}

The IPCHA protocol is implemented and deployed as an internal
research prototype within the nyxCore platform, a private multi-tenant
dashboard system.  The current deployment serves as an internal pilot
for the author's own workflows; it has not undergone independent
security audit or compliance certification, and the production-readiness
gaps identified in Section~\ref{sec:limitations} (audit log
immutability, authority document integrity, bypass dynamics) remain
unresolved.  We use the term ``deployed'' to indicate functional
integration, not compliance-ready production status.

The IPCHA verification sidecar, NLI microservice, evaluation corpus
(110~cases), and scoring scripts are available at
\url{https://github.com/mrwind-up-bird/ipcha-mistabra} under the
Apache~2.0 license.  The nyxCore platform, which provides the
dashboard and workflow integration, is a proprietary system not
required for reproducing the evaluation results.  The ONNX-exported
DeBERTa model is publicly available on HuggingFace
(\texttt{MoritzLaurer/DeBERTa-v3-base-mnli-fever-anli}).

% ================================================================
\bibliographystyle{unsrtnat}
\bibliography{references}

\end{document}
